\chapter*{Conclusion générale}
\addcontentsline{toc}{chapter}{Conclusion générale}

Ce projet de fin d'études avait pour objectif de concevoir et de réaliser une solution d'aide à la migration d'une application Angular monolithique vers une architecture micro-frontend. Cette problématique répond à un besoin réel dans les projets web modernes, où la croissance du code, la complexité fonctionnelle et les difficultés de maintenance rendent progressivement les architectures monolithiques moins adaptées. Le recours aux micro-frontends permet de mieux séparer les responsabilités, de favoriser l'autonomie des équipes et de rendre l'évolution applicative plus maîtrisable.

Pour répondre à cet objectif, le travail a d'abord consisté à étudier le contexte général du projet, les limites d'une application frontend monolithique et les principes d'une architecture micro-frontend. Cette étude a permis de comprendre les enjeux de la migration, notamment la séparation fonctionnelle, la conservation du comportement existant, la gestion des dépendances, la configuration de Module Federation et la nécessité de produire une documentation claire du résultat généré.

La solution proposée repose sur un pipeline de migration structuré et progressif. Ce pipeline analyse le projet Angular source, produit un plan de migration, génère une architecture cible composée d'un Shell et de plusieurs Remotes, migre les éléments applicatifs puis applique des validations techniques. L'utilisation d'agents spécialisés permet de séparer les responsabilités : analyse, conception cible, génération structurelle, migration des services, migration des composants, routage, configuration de fédération et documentation. Cette organisation rend la solution plus lisible, plus maintenable et plus facilement extensible.

Les différentes releases et les sprints réalisés ont permis de construire progressivement la solution. La première release a posé les bases du pipeline, de l'analyse du projet source et de la génération de la structure micro-frontend. La deuxième release a approfondi la migration détaillée des services, de la gestion d'état, des routes, des composants, des templates et des styles, avant de finaliser l'intégration avec Module Federation. Enfin, la partie réalisation a montré l'utilisation concrète du migrateur, la génération du dossier \texttt{output}, le lancement du nouveau projet et la production de la documentation finale.

Le résultat obtenu constitue une base fonctionnelle pour automatiser une partie importante d'une migration Angular vers une architecture micro-frontend. La solution ne se limite pas à copier des fichiers : elle organise la migration, conserve des rapports intermédiaires, applique des validations et fournit une documentation destinée au développeur. Ces éléments améliorent la traçabilité et facilitent la reprise du projet généré après l'exécution du pipeline.

Cependant, certaines limites subsistent. Les applications Angular peuvent présenter des structures très différentes, des conventions internes, des dépendances spécifiques ou des comportements difficiles à interpréter automatiquement. Dans ces cas, l'intervention humaine reste nécessaire pour valider certains choix de découpage, corriger des cas particuliers ou compléter les tests fonctionnels. De plus, la réussite du build ne garantit pas à elle seule que toutes les fonctionnalités métier se comportent exactement comme dans le projet initial.

Plusieurs perspectives d'amélioration peuvent être envisagées. Il serait possible d'enrichir l'analyse statique afin de mieux comprendre les dépendances complexes, d'améliorer les stratégies de découpage fonctionnel, d'ajouter des tests automatisés sur le projet généré et de renforcer l'interface graphique avec des visualisations plus détaillées. L'outil pourrait également être étendu à d'autres frameworks frontend ou à d'autres stratégies de décomposition applicative.

En conclusion, ce projet a permis de mettre en place une solution structurée, documentée et extensible pour accompagner la migration d'une application Angular monolithique vers une architecture micro-frontend. Il représente une contribution pratique à la modernisation des applications frontend et constitue une base solide pour des travaux futurs autour de l'automatisation, de la validation et de la gouvernance des migrations applicatives.
